% suplementario.tex - Material Suplementario (validacion estatica BSF)
% Compilar: latexmk -xelatex suplementario.tex
\documentclass[11pt,a4paper]{article}

\usepackage{fontspec}
\setmainfont{TeX Gyre Termes}
\setsansfont{TeX Gyre Heros}
\setmonofont{Latin Modern Mono Light}

\usepackage[spanish,es-noquoting,es-nodecimaldot]{babel}
\usepackage{amsmath}
\usepackage[margin=2.5cm]{geometry}
\usepackage{graphicx}
\usepackage{booktabs}
\usepackage{siunitx}
\usepackage{caption}
\usepackage{tikz}

\usepackage[backend=biber,style=authoryear,sorting=nyt]{biblatex}
\addbibresource{bibliografia/bibliografia.bib}

\usepackage[hidelinks]{hyperref}

% Numeracion S independiente
\renewcommand{\thesection}{S\arabic{section}}
\renewcommand{\theequation}{S\arabic{equation}}
\renewcommand{\thefigure}{S\arabic{figure}}
\renewcommand{\thetable}{S\arabic{table}}

\title{Material Suplementario\\[6pt]
{\large Modelado y validación experimental de emisiones de CO$_2$
  con separación de fuentes larvaria y microbiana en bioconversión
  estática con \textit{Hermetia illucens}}}
\author{}
\date{}

\begin{document}
\maketitle

\noindent Este Material Suplementario recoge el detalle matemático y de
protocolo que complementa al texto principal.

\tableofcontents
\bigskip

\section{Balance de gas en cámara estática}
\label{sec:camara}

El tercer componente conecta las fuentes con la medición: durante un cierre, el CO$_2$ se acumula en el volumen de aire de la cámara y el sensor registra la pendiente de concentración (Figura \ref{fig:camara}).

\begin{figure}[htbp]
\centering
\begin{tikzpicture}[
  comp/.style={draw, rounded corners=2mm, minimum width=2.6cm, minimum height=1.1cm, align=center, fill=gray!8},
  inst/.style={draw, rounded corners=2mm, minimum width=2.6cm, minimum height=1.1cm, align=center},
  flujo/.style={->, >=latex, thick, shorten >=3pt, shorten <=3pt},
  etiqueta/.style={fill=white, inner sep=1.5pt, font=\footnotesize}]
\node[comp] (larvas) at (0,1.8) {Larvas\\ $N$};
\node[comp] (lecho) at (0,-1.8) {Lecho\\ $D_M$, $W$};
\node[comp] (cámara) at (6.4,0) {Cámara cerrada\\ $V_{\mathrm{aire}}$};
\node[inst] (sensor) at (12.2,0) {Sensor\\ pendiente $dc/dt$};
\draw[flujo] (larvas.east) -- node[etiqueta, sloped, above]{$N\,r_{CO_2}^{\mathrm{lar}}$} (cámara.170);
\draw[flujo] (lecho.east) -- node[etiqueta, sloped, below]{$r_C^{\mathrm{mic}}$} (cámara.190);
\draw[flujo] (cámara) -- node[etiqueta, above=2pt]{$R_{CO_2}/V_{\mathrm{aire}}$} (sensor);
\end{tikzpicture}
\caption{Balance de gas en la cámara estática, sobre el método estándar \parencite{hutchinsonMethodsSoilAnalysis1992}. Las fuentes larvaria y microbiana acumulan gas en el volumen de aire durante un cierre (ec. \ref{eq:camara}); el sensor registra la pendiente, llevada a ppm/min con la ley de gas ideal (ec. \ref{eq:ppm}).}
\label{fig:camara}
\end{figure}

\paragraph{Acumulación de gas durante el cierre de la cámara.}
El método de cámara estática estima la tasa a partir de la acumulación lineal del gas en un volumen conocido \parencite{hutchinsonMethodsSoilAnalysis1992}; con las fuentes del modelo, el balance molar sobre el aire de la cámara es
\begin{equation}
\frac{dc_{CO_2}}{dt} = \frac{R_{CO_2}}{V_{\mathrm{aire}}},
\qquad
R_{CO_2} = \frac{N\,r_{CO_2}^{\mathrm{lar}} + 1000\,r_C^{\mathrm{mic}}}{44000},
\label{eq:camara}
\end{equation}
donde $c_{CO_2}$ es la concentración molar del aire de la cámara (mol/L), $V_{\mathrm{aire}}$ el volumen de aire efectivo (L) y $R_{CO_2}$ la fuente total (mol/d): el término larvario pasa de mg a g y ambos a mol con la masa molar del CO$_2$ (44~g/mol).

\paragraph{Escala del sensor (ppm/min).}
La comparación con el sensor se hace en la escala propia del método \parencite{hutchinsonMethodsSoilAnalysis1992}, convirtiendo la acumulación molar a ppm/min con la ley de gas ideal \parencite{atkinsPhysicalChemistry2014}:
\begin{equation}
r_{\mathrm{ppm}} = \frac{R_{CO_2}}{V_{\mathrm{aire}}}\,\frac{R_g\,T}{P}\,\frac{10^6}{1440}
\quad [\mathrm{ppm/min}],
\label{eq:ppm}
\end{equation}
donde $R_g = 8.314$~J/(mol$\cdot$K) es la constante de los gases, y $T$ y $P$ la temperatura y la presión de la cámara: el factor $R_g T/P$ es el volumen molar (L/mol), $10^6$ convierte la fracción molar a ppm y 1440 pasa de días a minutos.

\section{Extensiones del núcleo larvario}
\label{sec:nucleoext}

\paragraph{Asimilación: extensiones al régimen por lotes.}
A la expresión de la tasa de asimilación del modelo de referencia se le aplican dos extensiones que la adaptan al régimen por lotes:
\begin{equation}
a(B, t, D_M) = a_{\mathrm{E}}(B, t)\; f\; \frac{D_M}{D_M + K_s},
\qquad r_A = a\,B.
\label{eq:asimilacion}
\end{equation}
El factor de calidad de la dieta $f$ ($0 < f \le 1$) escala la asimilación cuando la dieta es menos nutritiva que la de referencia ($f = 1$ en D1; en D4 se calibra). El término $D_M/(D_M + K_s)$ limita la asimilación por disponibilidad de alimento: $K_s$, la constante de semisaturación, se fija en 2~g, pequeña frente a la carga agregada del lecho ($D_M(0) = 647.3$~g en D1 y $602.3$~g en D4), de modo que la asimilación opera en saturación durante todo el ensayo (factor $\geq 0.99$), sin parámetros adicionales calibrados. El flujo asimilado total sigue siendo proporcional a la estructura, $r_A = a\,B$, como en el modelo de referencia.

\paragraph{Mantenimiento: suelo de prepupa.}
En el modelo de referencia el mantenimiento es proporcional a la estructura, $r_{Cm} = m\,B$ \parencite{eriksenDynamicModellingFeed2022}. Como extensión se añade un suelo para la prepupa: al detenerse la asimilación en prepupa (instar 7), el mantenimiento cae a una fracción $m_{\mathrm{suelo}}$ del valor basal, anclada a la medición de larvas en ayuno del día 17 (resultados del texto principal):
\begin{equation}
r_{Cm} = m\,B\,\big[(1 - \varepsilon_{\mathrm{prep}}) + m_{\mathrm{suelo}}\,\varepsilon_{\mathrm{prep}}\big],
\qquad
\varepsilon_{\mathrm{prep}} =
\begin{cases}
1, & B \ge B_{\max}(t) \quad \text{(prepupa)},\\
0, & \text{en otro caso},
\end{cases}
\label{eq:mantenimiento}
\end{equation}

\paragraph{Interpretación del umbral de transición.}
El umbral de transición está escrito en términos de $B_{\max}$; en este trabajo se interpretó como la $B_{\max}(t)$ vigente en cada instante, de modo que el decrecimiento del techo induce efectivamente el cierre del ciclo larvario.

\section{Cinética microbiana del lecho}
\label{sec:micdetalle}

\paragraph{Constante efectiva de degradación.}
La constante efectiva de degradación $k$ (d$^{-1}$) combina tres elementos clásicos de la cinética de descomposición de residuos \parencite{roels1980,shulerbioprocess2017}: una corrección térmica tipo van 't Hoff, una respuesta de saturación a la humedad tipo Monod y un tope numérico:
\begin{equation}
k = \min\!\left(k_{\mathrm{ref}}\,Q_{10}^{(T - T_{\mathrm{ref}})/10}\,
\frac{\theta_s}{\theta_s + K_{\theta}},\; k_{\max}\right).
\label{eq:kmic}
\end{equation}
El primer término, $k_{\mathrm{ref}}$, es la tasa específica a la temperatura de referencia y es el único parámetro cinético calibrado. El segundo, $Q_{10}^{(T - T_{\mathrm{ref}})/10}$, es la corrección térmica: la tasa se multiplica por $Q_{10}$ cada 10~°C de aumento sobre $T_{\mathrm{ref}}$; se toman los valores nominales estándar $Q_{10} = 2$ y $T_{\mathrm{ref}} = 25$~°C \parencite{roels1980,shulerbioprocess2017}. El tercero, $\theta_s/(\theta_s + K_{\theta})$, es una forma de Monod sobre la humedad con semisaturación $K_{\theta} = 0.15$ (fija, no calibrada). Finalmente, el tope $k_{\max} = 1.0$~d$^{-1}$ (fijo) se incluye solo por seguridad numérica y nunca resulta activo.

\paragraph{Crecimiento de la biomasa microbiana.}
Al degradar el sustrato, los microorganismos incorporan una fracción del carbono a biomasa nueva: el crecimiento microbiano es proporcional a la tasa de degradación:
\begin{equation}
\frac{dX}{dt} = Y_X\,r_{\mathrm{deg}},
\qquad X(0) = X_0,
\label{eq:xmic}
\end{equation}
donde $X$ es la biomasa microbiana del lecho (g) e $Y_X$ el rendimiento de biomasa (g de biomasa por g de sustrato degradado), concepto estándar de la cinética de cultivos \parencite{roels1980,shulerbioprocess2017}.

\paragraph{Trayectoria de humedad del ensayo y efecto neto sobre $k$.}
En este ensayo $\theta_s$ partió de $0.14$ (D1) y $0.20$ (D4) y subió con el aporte de agua metabólica, de modo que el factor de humedad operó entre el $48$ y el $62$~\% de la saturación: la humedad limitó la degradación al inicio y su efecto se atenuó hacia el final; con los valores calibrados, $k \approx 0.24$--$0.28$~d$^{-1}$ según la dieta, lejos del tope $k_{\max}$.

\paragraph{Comprobación de la representación de la segunda carga.}
Representar la segunda carga de alimento (250~g al día~5) como un pulso explícito de sustrato seco y agua en $t = 4$~d, en lugar de la carga agregada al montaje, desplaza los parámetros estructurales calibrados menos de un 5~\% (solo $\tau$ se desplaza alrededor de un día, coherente con su débil identificabilidad individual, sec.~3.4 del texto principal), no modifica $k_{\mathrm{ref}}$ ni $Y_{CO_2}$ y deja los tiempos de cierre emergentes prácticamente sin cambio (11.4 a 11.6~d en D1; 12.5~d en D4).

\section{Cotas de búsqueda del ajuste de parámetros}
\label{sec:cotas}

Las cotas de búsqueda no son arbitrarias: cada intervalo encierra el rango físicamente plausible de su parámetro. Las de $k_{\mathrm{ref}}$ abarcan tres órdenes de magnitud, desde una degradación despreciable ($10^{-3}$~d$^{-1}$) hasta la cota cinética del propio modelo ($k_{\max} = 1$~d$^{-1}$). Las de $Y_{CO_2}$ admiten desde una mineralización casi nula hasta 2~g de CO$_2$ por gramo de sustrato, valor superior a 1 posible porque la oxidación incorpora oxígeno del aire (la oxidación completa de carbohidratos rinde $\approx 1.5$~g~g$^{-1}$). El intervalo de $B_0$ (0.05 a 2~mg) cubre la estructura esperable de una larva de 4.5~d de edad según las trayectorias del modelo \parencite{eriksenDynamicModellingFeed2022}. El intervalo de $f$ permite desde una dieta cinco veces peor que la de referencia (0.2) hasta una ligeramente mejor (1.2). Las cotas de $a_{\max}$ (1.1 a 1.5~d$^{-1}$) son el rango reportado por \textcite{eriksenDynamicModellingFeed2022} para dietas estándar; las de $B_{\max,0}$ (20 a 80~mg) abarcan desde larvas de talla pequeña hasta más allá de la cepa de 65~mg de \textcite{eriksenDynamicModellingFeed2022}; las de $\rho$ (0.5 a 15~mg~d$^{-1}$) parten del valor de referencia (1~mg~d$^{-1}$) y admiten un decrecimiento mucho más rápido a la temperatura del ensayo; las de $\delta_{DW,f}$ (0.28 a 0.65) van desde el valor medido al montaje hasta por encima del rango de conversión a peso seco (0.20 a 0.45) \parencite{eriksenDynamicModellingFeed2022}, para no forzar la solución; y las de $\tau$ (0 a 6~d) acotan la anticipación de la rampa de materia seca respecto al cierre emergente.

\section{Tasas netas de CO$_2$ por dieta y día de medición}
\label{sec:tasasdia}

\begin{table}[htbp]
\centering
\small
\caption{Tasas netas de CO$_2$ por dieta y día de medición (ppm/min): control de alimento solo observado, tasa neta observada (tratamiento menos control) y componente larvaria del modelo en el ajuste conjunto. ($^{*}$) cierre saturado: la tasa observada es una cota inferior. La desviación es (modelo $-$ observado)/observado; para las cotas, un valor negativo indica cota no alcanzada. Los días del experimento corresponden a los días desde el montaje más el desfase $\Delta = 2$~d (sec. de solución numérica del texto principal).}
\label{tab:tasas}
\begin{tabular}{ll rrr r}
\toprule
Dieta & Día & Control obs. & Neta obs. & Larvaria mod. & Desv. (\%) \\
\midrule
D1 & 9  & 26.5 & 267.8$^{*}$ & 134.2 & $-49.9$ \\
   & 11 & 34.8 & 124.3$^{*}$ & 105.2 & $-15.4$ \\
   & 13 & 40.2 & 102.8        & 77.8  & $-24.3$ \\
\midrule
D4 & 9  & 25.8 & 88.6         & 134.2 & $+51.4$ \\
   & 11 & 27.2 & 165.4$^{*}$ & 111.1 & $-32.8$ \\
   & 13 & 34.6 & 21.5         & 74.0  & $+244$ \\
   & 17 & 54.1 & 0.0          & 9.6   & -- \\
\bottomrule
\end{tabular}
\end{table}

\section{Sensibilidad de la calibración: tabla completa, recuperación por escenario y bootstrap}
\label{sec:senstab}

\begin{table}[htbp]
\centering
\footnotesize
\caption{Sensibilidad de la calibración conjunta a los parámetros fijos: parámetros recuperados, calidad de ajuste y cierre emergente en cada escenario.}
\label{tab:sensibilidad}
\begin{tabular}{l ccccc cc cc cc}
\toprule
& \multicolumn{5}{c}{Parámetros recuperados} & \multicolumn{2}{c}{RMSE gas} & \multicolumn{2}{c}{RMSE biomasa} & \multicolumn{2}{c}{$t_{\mathrm{stop}}$} \\
\cmidrule(lr){2-6} \cmidrule(lr){7-8} \cmidrule(lr){9-10} \cmidrule(lr){11-12}
Escenario & $B_0$ & $f$ & $a_{\max}$ & $B_{\max,0}$ & $\rho$ & D1 & D4 & D1 & D4 & D1 & D4 \\
\midrule
Base & 1.10 & 0.802 & 1.50 & 48.9 & 2.00 & 25.0 & 40.5 & 24.5 & 12.1 & 11.4 & 12.5 \\
$m \times 0.8$ & 1.09 & 0.806 & 1.46 & 55.8 & 2.50 & 31.7 & 41.7 & 20.8 & 12.0 & 11.1 & 12.1 \\
$m \times 1.2$ & 1.26 & 0.745 & 1.44 & 51.1 & 2.66 & 7.1 & 50.0 & 26.3 & 18.0 & 11.1 & 12.5 \\
$Y_B \times 0.8$ & 1.03 & 0.775 & 1.50 & 51.7 & 2.12 & 20.2 & 42.3 & 18.4 & 11.4 & 11.3 & 12.5 \\
$Y_B \times 1.2$ & 1.15 & 0.813 & 1.49 & 50.0 & 2.23 & 23.9 & 45.2 & 27.8 & 14.0 & 11.3 & 12.3 \\
$N \times 0.8$ & 1.02 & 0.786 & 1.50 & 61.5 & 2.89 & 25.1 & 41.3 & 19.3 & 14.6 & 11.1 & 12.3 \\
$N \times 1.2$ & 0.92 & 0.890 & 1.50 & 41.9 & 1.39 & 22.7 & 44.2 & 36.0 & 11.3 & 11.9 & 12.6 \\
$t_{pm} \times 0.8$ & 1.27 & 0.780 & 1.41 & 52.7 & 1.88 & 22.6 & 42.7 & 24.8 & 13.9 & 11.1 & 12.8 \\
$t_{pm} \times 1.2$ & 1.14 & 0.767 & 1.50 & 51.3 & 6.15 & 19.8 & 44.2 & 21.6 & 13.3 & 11.1 & 11.5 \\
\bottomrule
\end{tabular}
\end{table}

\noindent En el Cuadro~\ref{tab:sensibilidad}, cada fila perturba un parámetro fijo ($\times 0.8$ o $\times 1.2$) y recalibra los ocho grados de libertad; la fila base es la calibración nominal. Unidades: $B_0$ y $B_{\max,0}$ en mg, $a_{\max}$ en d$^{-1}$, $\rho$ en mg~d$^{-1}$, RMSE de gas en ppm/min, RMSE de biomasa en mg de peso húmedo y $t_{\mathrm{stop}}$ en días desde el montaje. En todas las filas la cota saturada del día 9 de D1 y la del día 11 de D4 quedan sin alcanzar (ver texto).

\paragraph{Recuperación de los ocho grados de libertad por escenario.}
Los ocho grados de libertad se recuperan de forma estable en los nueve escenarios: $B_0$ entre 0.92 y 1.27~mg, $f$ entre 0.75 y 0.89 y $B_{\max,0}$ entre 41.9 y 61.5~mg, con una calidad de ajuste apenas alterada (RMSE de biomasa de a lo sumo 36~mg; RMSE de gas de 7.1 a 31.7~ppm/min en D1 y de 40.5 a 50.0 en D4, alrededor de los valores base). $a_{\max}$ permanece en su cota superior en el caso base y en cuatro de los ocho escenarios perturbados (en los demás, entre 1.41 y 1.49~d$^{-1}$): el dato sigue pidiendo una asimilación temprana mayor que la permitida, independientemente de los parámetros fijos.

\paragraph{Bootstrap de residuos del ajuste base.}
Como complemento, la incertidumbre de muestreo de los ocho grados de libertad se evaluó con un bootstrap de residuos: 200 réplicas en las que los residuos crudos del ajuste base se remuestrearon con reemplazo por dieta y por bloque (gas y biomasa), se sumaron a los valores ajustados y se repitió la calibración completa; los puntos saturados conservaron su tratamiento de cota inferior y el módulo microbiano se mantuvo fijo. Los intervalos resultantes son conservadores (los residuos remuestreados incluyen la discrepancia estructural de los cierres saturados). Los percentiles 2.5 y 97.5 confirman el cuadro de identificabilidad: $f$ queda entre $0.54$ y $0.91$ (mediana $0.73$) y $B_0$ entre $0.91$ y $2.0$~mg, mientras $a_{\max}$ ($1.1$--$1.5$~d$^{-1}$), $\rho$ ($0.5$--$5.1$~mg~d$^{-1}$), $\delta_{DW,f}$ ($0.35$--$0.64$) y los $\tau$ ($0.1$--$6.0$~d en D1 y $2.1$--$6.0$~d en D4) cubren todo o casi todo su rango permitido. El cierre emergente se mantiene como la cantidad reportable: mediana de $11.4$~d en D1 (percentil 2.5: $11.1$~d) y de $13.1$~d en D4 (percentil 2.5: $11.7$~d), con el ciclo cerrando dentro de la ventana de medición en el 95~\% (D1) y el 85~\% (D4) de las réplicas.

\section{Validación externa con los datos de Schøn et al.: tabla por dieta}
\label{sec:schntab}

\begin{table}[htbp]
\centering
\small
\caption{Validación externa con los datos de Schøn et al.: factor de calidad ajustado y comparación de volumen total, día del pico y masa final por dieta.}
\label{cua:schn}
\begin{tabular}{l c rr rr rr}
\toprule
& & \multicolumn{2}{c}{VCO$_2$ total} & \multicolumn{2}{c}{Día del pico} & \multicolumn{2}{c}{Masa final} \\
\cmidrule(lr){3-4} \cmidrule(lr){5-6} \cmidrule(lr){7-8}
Dieta & $f$ & obs. & mod. & obs. & mod. & obs. & mod. \\
\midrule
Alimento de pollo & 0.907 & 19.6 & 15.0 & 5.7 & 5.2 & 43.2 & 51.6 \\
P:C 1:3           & 0.883 & 16.9 & 14.7 & 6.0 & 5.5 & 37.7 & 50.3 \\
P:C 1:5           & 0.878 & 16.2 & 14.7 & 6.0 & 5.5 & 33.5 & 50.0 \\
P:C 1:9           & 0.733 & 10.8 & 12.6 & 5.5 & 7.0 & 21.0 & 40.2 \\
\bottomrule
\end{tabular}
\end{table}

\noindent Unidades del Cuadro~\ref{cua:schn}: volumen total de CO$_2$ en litros por lote (600 larvas, 80~g de materia seca de alimento, 7~d), día del pico desde el inicio de su ensayo y masa final en mg de materia seca por larva.

\paragraph{Detalle por dieta: volúmenes, picos y masas finales.}
Los volúmenes totales de CO$_2$ quedan dentro del $\pm 25$~\% en las cuatro dietas (figura de validación del texto principal): el modelo subestima entre un 9 y un 23~\% las tres dietas de mejor calidad y sobreestima un 17~\% la más pobre. El día del pico de emisión se reproduce dentro de medio día en las tres primeras dietas; en la 1:9 el pico del modelo llega 1.5~d más tarde que el observado, en el borde de la ventana de medición (7~d).

Las masas finales quedan entre 1.19 y 1.90 veces las observadas, con la sobreestimación creciendo al empeorar la dieta: en la dieta 1:9 las larvas observadas terminan con menos de la mitad de la masa de las alimentadas con alimento de pollo (21.0 frente a 43.2~mg de materia seca), un retardo del desarrollo que el núcleo solo puede representar reduciendo la asimilación vía $f$. Además, en ese estudio el flujo medido es el total larvas--lecho (sin control de alimento solo): la prueba también extrapola el módulo microbiano a otra temperatura y otro sustrato, y que aun así los totales queden dentro del $\pm 25$~\% y los picos dentro de medio día para las dietas de calidad media y alta respalda la transferibilidad de la parametrización conjunta.

\section{Consolidación de las validaciones}
\label{sec:consol}

% Generada por simulacion/consolidacion_validaciones.py (Fase B)
\begin{table}[htbp]
  \centering
  \caption{Consolidaci\'on de las validaciones del modelo (bloques B0--B6). B0--B4 corresponden al ajuste y a los chequeos con datos propios; B5--B6 son validaciones externas por transferencia de par\'ametros. Los RMSE de B1 est\'an en ppm/min y los de B2 en mg de peso h\'umedo por larva.}
  \label{tab:consolidacion_validaciones}
  \small
  \resizebox{\textwidth}{!}{%
  \begin{tabular}{l l l c c c l}
    \hline
    Bloque & Fuente & M\'etrica & Modelo & Observado & Criterio & Resultado \\
    \hline
    B0 & Propia (pesadas d0) & Masa inicial $B_0 + L_0$ (mg MS/larva) & 1.43 & 1.3 / 0.9 (D1/D4) & orden de magnitud & CUMPLE \\
    B1 & Propia (camara estatica) & RMSE flujo CO$_2$ D1 (ppm/min) & 25.1 & --- & --- & Ajuste \\
    B1 & Propia (camara estatica) & RMSE flujo CO$_2$ D4 (ppm/min) & 39.3 ($R^2 = -0.08$) & --- & --- & Limitacion declarada \\
    B2 & Propia & RMSE biomasa D1 / D4 (mg PH) & 25.1 / 15.2 & --- & --- & Ajuste \\
    B2 & Propia & Masa larval d14 D1 / D4 (mg PH) & 169 / 144 & 173 / 135 & $\pm$10\,\% & CUMPLE \\
    B3 & Propia (bromatologia) & Lipidos d14 D1 / D4 (\% MS) & 34 / 28 & 33.4 / 26.6 (EE) & punto a punto & CUMPLE \\
    B4 & Propia (ec. 22 Eriksen) & Fraccion MS $\delta_{DW}$ & 0.39 & 0.33--0.45 & dentro del rango & CUMPLE \\
    B5 & Ermolaev et al. & Tasas CO$_2$ (3 pts) + acumulado d14 (g kg$^{-1}$) & 90.8 & $92.4 \pm 14.5$ & residuos $< 1\sigma$ & CUMPLE \\
    B5 & Ermolaev et al. & Masa larval final (mg) & 111.1 & 101.9 & factor 2 & CUMPLE ($+9$\,\%) \\
    B5 & Ermolaev et al. & $t_{pm}$ (d) / $B_{max}$ (mg) recuperados & 13.0 / 94.8 & 13 / 90 (Eriksen) & consistencia & CUMPLE \\
    B6 & Sch\o{}n et al. & $f$ por dieta (REF, 1:3, 1:5, 1:9) & 0.899, 0.874, 0.870, 0.725 & orden P:C & coherente & CUMPLE \\
    B6 & Sch\o{}n et al. & VCO$_2$ total (L) & 15.0, 14.7, 14.7, 12.6 & 19.6, 16.9, 16.2, 10.8 & $\pm$25\,\% & CUMPLE \\
    B6 & Sch\o{}n et al. & Masa final (mg MS) & 51.5, 50.0, 49.8, 40.0 & 43.2, 37.7, 33.5, 21.0 & factor 2 & CUMPLE \\
    \hline
  \end{tabular}}%
\end{table}


\section{Detalles del montaje experimental}
\label{sec:montaje}

Las mediciones de gas se realizaron en cámaras plásticas rígidas de
geometría troncopiramidal rectangular (boca de $37 \times 26$~cm, base de
$32 \times 21$~cm y altura de $17$~cm; volumen geométrico aproximado de
13.8~L). Durante cada cierre la cámara se mantuvo sellada con su tapa
reforzada con silicona, de modo que la concentración interior de CO$_2$ solo
cambia por la producción del contenido. El volumen de aire efectivo
$V_{\mathrm{aire}}$ ---el volumen geométrico menos el ocupado por el
recipiente interno y su contenido--- entró como parámetro medido en el cálculo
de tasas.

El sensor de CO$_2$ (MH-410D) se montó empotrado en una de las caras
laterales mayores, con el elemento sensible hacia el interior y por encima
del nivel del sustrato. Antes del experimento se verificó en laboratorio
contra un cromatógrafo de gases con fuente certificada, en el rango de
medición del experimento (0--6000~ppm).

El sustrato y las larvas se alojaron en un recipiente plástico de
$15 \times 12 \times 5$~cm con tapa aireada por malla, colocado dentro de la
cámara; se supuso equilibrio rápido de gases entre el recipiente interno y el
aire de la cámara a través de la malla. El mismo tipo de montaje se usó en
tratamientos y controles, de modo que $V_{\mathrm{aire}}$ fue comparable
entre ellos.

Las cargas de alimento se dosificaron tal como se recibió, sin ajuste de humedad; con la materia seca del análisis bromatológico (86.3~\% en D1 y 80.3~\% en D4), la primera carga de 500~g equivale a 431.5 y 401.5~g de materia seca en D1 y D4, respectivamente, y la segunda de 250~g a 215.8 y 200.8~g.

\section{Validación externa con los datos de Ermolaev et al.: tabla día a día}
\label{sec:ermoltab}

\begin{table}[htbp]
\centering
\small
\caption{Validación externa con los datos de Ermolaev et al.: tasas puntuales y emisión acumulada, observadas y modeladas.}
\label{cua:ermolaev}
\begin{tabular}{c cc cc}
\toprule
& \multicolumn{2}{c}{Tasa} & \multicolumn{2}{c}{Acumulada} \\
\cmidrule(lr){2-3} \cmidrule(lr){4-5}
Día & obs. & mod. & obs. & mod. \\
\midrule
2  & $2.53 \pm 1.1$ & 2.52 & 9.1  & 3.6 \\
7  & $8.04 \pm 1.6$ & 7.09 & 37.3 & 41.2 \\
11 & $5.58 \pm 1.5$ & 6.08 & 68.6 & 73.1 \\
14 & --- & --- & $92.4 \pm 14.5$ & 90.7 \\
\bottomrule
\end{tabular}
\end{table}

\noindent Unidades del Cuadro~\ref{cua:ermolaev}: tasas en mg~min$^{-1}$ para el lote completo del reactor y acumulada en g~kg$^{-1}$ de alimento en peso fresco; las incertidumbres observadas son una desviación estándar (tasas) o el intervalo reportado (acumulada).

Quedan a la vista dos desacuerdos menores: el modelo subestima la emisión acumulada del arranque (día 2: 3.6 frente a 9.1~g~kg$^{-1}$) aunque reproduce la tasa puntual de ese día casi exactamente, y su tasa máxima ocurre hacia el día 5.6, alrededor de un día y medio antes del pico observado. Ambos rasgos apuntan a una subestimación de la respiración en los primeros días, la misma discrepancia de fase temprana encontrada con los datos propios (sec.~3.4 del texto principal).

\section{Parámetros fijos y de literatura del modelo}
\label{sec:paramfijos}

\begin{table}[htbp]
\centering
\caption{Parámetros fijos y tomados de la literatura: símbolo, valor y unidad. La procedencia de cada valor se detalla en el texto principal.}
\label{tab:paramfijos}
\begin{tabular}{llll}
\toprule
Parámetro & Símbolo & Valor & Unidad \\
\midrule
Inicio del decrecimiento de $B_{\max}$ (fijo) & $t_{pm}$ & 8.5 (13 desde eclosión) & d \\
Divisor de instares (fijo) & $\delta_I$ & 5 & -- \\
Exponente de asimilación (fijo) & $\alpha$ & 1.0 & -- \\
Exponente de crecimiento (fijo) & $\beta$ & 2.0 & -- \\
Costo de crecimiento & $Y_B$ & 0.44 & -- \\
Costo de depósito de lípidos & $Y_L$ & 0.42 & -- \\
Mantenimiento & $m$ & 0.08 & d$^{-1}$ \\
Suelo de mantenimiento & $m_{\mathrm{suelo}}$ & 0.13 & -- \\
Factor $Q_{10}$ microbiano & $Q_{10}$ & 2.0 & -- \\
Semisaturación de humedad (fija) & $K_{\theta}$ & 0.15 & -- \\
Cota cinética microbiana (fija) & $k_{\max}$ & 1.0 & d$^{-1}$ \\
Biomasa microbiana inicial (fija) & $X_0$ & 1 & g \\
Rendimiento de biomasa microbiana (fijo) & $Y_X$ & 0.4 & g/g \\
Semisaturación de alimento (fija) & $K_s$ & 2 & g \\
Semisaturación de sustrato microbiano (fija) & $K_{sm}$ & 5 & g \\
\bottomrule
\end{tabular}
\end{table}

\section{Análisis de incertidumbre de las tasas de CO$_2$}
\label{sec:incertidumbre_si}

Las fuentes de incertidumbre consideradas fueron: (i) la exactitud instrumental del sensor de CO$_2$ (MH-410D, $\pm(50$ ppm $+\ 5$ \% de la lectura) según el fabricante \parencite{winsenMH410D2022}, verificado contra cromatógrafo (sec.~2.1 del texto principal)), propagada a la pendiente de cada cierre como un error relativo de $5$ \% $+\ 50/\Delta C$, con $\Delta C$ el incremento de concentración del cierre (entre 6 y 12 \% para las tasas de este trabajo); (ii) la calidad del ajuste lineal de cada curva, controlada con el filtro $R^2 \geq 0.5$ y el tratamiento de los cierres saturados como cotas inferiores; (iii) la dispersión entre los cierres consecutivos del mismo lote; y (iv) la incertidumbre del volumen de aire y de las condiciones de la cámara (menor al 2 \%). La incertidumbre de la tasa neta se propagó como $\sigma_{\mathrm{neto}} = \sqrt{\sigma_{\mathrm{trat}}^2 + \sigma_{\mathrm{ctrl}}^2}$, suponiendo errores independientes entre tratamiento y control del mismo día. Para la tasa molar $R = f(m, V_{\mathrm{aire}}, T)$, con $m$ la pendiente medida, la ley de propagación para errores independientes da
\begin{equation}
\left(\frac{u_R}{R}\right)^2 = \left(\frac{u_m}{m}\right)^2 + \left(\frac{u_V}{V_{\mathrm{aire}}}\right)^2 + \left(\frac{u_T}{T}\right)^2,
\label{eq:propagacion}
\end{equation}
con $u_m/m \approx 0.06$ a $0.12$ (instrumental, sistemática: afecta todos los cierres en el mismo sentido y no se reduce al promediar), $u_V/V_{\mathrm{aire}} < 0.02$ y $u_T/T < 0.005$, lo que da una incertidumbre combinada de 6.5 a 12 \% por tasa. La componente aleatoria domina la dispersión punto a punto observada entre el modelo y los datos. El tiempo de respuesta del sensor ($T_{90} < 30$ s) es despreciable frente a la duración de los cierres.

\section{Validación externa con datos independientes: detalle de la transferencia}
\label{sec:validaciones_si}

Para evaluar la transferibilidad a condiciones distintas de las de calibración se emplearon los datos de \textcite{ermolaevGreenhouseGasEmissions2019}: 700 larvas de \textit{H. illucens} sobre residuos de alimentos (24.9~\% de materia seca) a 28~$^{\circ}$C durante 14 días, con alimentación en pulsos los días 0, 2, 7 y 11 (294~g de materia seca, unas cuatro veces la ración por larva de este trabajo). Las tasas puntuales (días 2, 7 y 11) y la emisión acumulada se digitalizaron de su figura~2, con error menor al 2~\% frente a los valores reportados en su texto y tablas. La simulación transfiere sin cambios la cinética larvaria y el módulo microbiano calibrados: el estado larvario inicial se tomó de la simulación propia a los cinco días de edad, los pulsos se representaron como adiciones discretas de sustrato seco y agua según la humedad reportada, y la temperatura se aplicó mediante la corrección $Q_{10}$ del módulo microbiano. Como el ensayo cubre el ciclo larvario completo (más allá de la ventana de medición propia), se dejaron libres los cuatro parámetros asociados a la dieta y al cierre del ciclo ---$f$, $t_{pm}$, $\rho$ y $B_{\max,0}$---, ajustados por mínimos cuadrados a las tres tasas puntuales y a la emisión acumulada ($92.4 \pm 14.5$~g~kg$^{-1}$), con residuos ponderados por la incertidumbre reportada; todos los demás parámetros permanecieron fijos en sus valores calibrados.

Como segunda prueba se emplearon los datos de \textcite{schonUsingGasExchange2026}: intercambio gaseoso de 600 larvas de 5--6 días (5.0~mg en peso húmedo) sobre cuatro dietas de composición proximal conocida (alimento de pollo y tres mezclas semisintéticas con proporciones proteína:carbohidrato de 1:3, 1:5 y 1:9; 19.9 a 7.0~\% de proteína), en lotes únicos de 80~g de materia seca y 261~g de agua, a 27~$^{\circ}$C durante 7 días, con respirometría de flujo abierto (el flujo medido es la producción total larvas--lecho; el control fue una cámara vacía, sin control de alimento solo). Las curvas horarias de consumo de O$_2$ y de relación de intercambio respiratorio (RER) se digitalizaron de su figura~2 y se combinó VCO$_2$ = VO$_2\,\times\,$RER, reescalando cada curva para que su integral coincida con los totales reportados en su tabla~2 (error final menor al 2~\%). El estado larvario inicial se tomó de la simulación propia en el instante en que la biomasa alcanza los 1.3~mg de materia seca de sus larvas iniciales (sus animales, de una colonia comercial, llegan más pequeños a igual edad). La transferencia fija de nuevo el núcleo DEB completo y el módulo microbiano calibrado, y deja libre únicamente $f$ por dieta, ajustado por mínimos cuadrados sobre las cuatro curvas horarias reconstruidas (29 puntos por dieta). Los valores de RER mayores que 1 (lipogénesis) son estructuralmente ajenos al modelo, que no representa O$_2$; la validación se limita por tanto al VCO$_2$.

\printbibliography

\end{document}
